\section*{Chapitre \ref{ch:m3:spot-markets} --- Marchés au comptant}

\begin{solution}{exo:m3:spot-markets:1}
$3\,000/60\,000 = 0.05$. Trois frais de 10 points de base forment une bande de $-3\ln(1 - 0.001) \approx 30$ points
de base de chaque côté, auxquels s'ajoutent les fourchettes acheteur-vendeur~: environ 0.04985 à 0.05015 avant prise en compte des fourchettes.
\end{solution}

\begin{solution}{exo:m3:spot-markets:2}
$60\,000/60\,030 = 0.9995$~: l'USDT à environ 99.95 cents rend les prix cohérents.
\end{solution}

\begin{solution}{exo:m3:spot-markets:3}
Les crypto-actifs circulent librement entre les plateformes, si bien que le cours ether--bitcoin reste aligné grâce aux seuls transferts de crypto-actifs~; le
cours du bitcoin en won face au dollar exige la jambe en monnaie fiduciaire, que les contrôles des capitaux et les règles de comptes restreignent.
\end{solution}

\begin{solution}{exo:m3:spot-markets:4}
$\Phi(-80/(6\,000\sqrt{30/525\,600})) = 3.9\%$ sur 30 minutes et 10.6~\% sur 60 minutes.
\end{solution}

\begin{solution}{exo:m3:spot-markets:5}
L'indice de référence compte 1.5 transaction non arrondie par transaction arrondie. Les transactions authentiques de la bourse représentent environ $12\% \times 2.5 =
30\%$ de ses transactions, si bien que le \omterm{def:m3:spot-markets:wash}{wash trading} représente environ 70~\%.
\end{solution}

\begin{solution}{exo:m3:spot-markets:6}
Les deux jambes s'exécutent en même temps, si bien que l'écart est verrouillé. L'inventaire lui-même subsiste~: son risque de prix (généralement couvert
sur une plateforme de contrats à terme), la défaillance des plateformes qui le détiennent, le coût et le délai de rééquilibrage lorsque les flux sont à sens unique,
et le capital immobilisé.
\end{solution}

\begin{solution}{exo:m3:spot-markets:7}
(triangulaire, croisé) = (1\,254, 2\,953), (794, 1\,902), (309, 793) et (26, 120). Chaque instantané offre 6
triangles (deux sens sur trois plateformes), mais 18 \omterm{def:m3:spot-markets:cross}{arbitrages entre plateformes} (trois paires, six couples de plateformes ordonnés), et un
transaction croisée paie deux frais contre trois pour un triangle.
\end{solution}

\begin{solution}{exo:m3:spot-markets:8}
Le volume déclaré peut relever en grande partie du \omterm{def:m3:spot-markets:wash}{wash trading}, et même un volume réel renseigne peu sur le coût d'un ordre~: comparer
la \omterm{def:m3:spot-markets:depth}{profondeur affichée} près du prix milieu, mesurée dans le temps, et le glissement de prix réalisé sur ses propres transactions.
\end{solution}

\begin{solution}{pb:m3:spot-markets:1}
\textbf{1.} $1.001 \times 1.0005$~: 1.0015 dollars, soit 0.15~\% de coûts. \textbf{2.} $1.65 \times 0.70 \times
\sqrt{60/525\,600} = 1.23\%$. \textbf{3.} $0.9995 \times 0.995 \times 0.98 = 97.46\%$. \textbf{4.} $1.0015
\times 1.0123/0.9746 - 1 = 4.03\%$. \textbf{5.} $1.15 \times 0.9746/1.0015 - 1 = 11.9\%$ par aller-retour.
\textbf{6.} 1.95~\%. \textbf{7.} 3.28~\%. \textbf{8.} Environ 5\,956~USD. \textbf{9.} Les bourses coréennes servaient les résidents
disposant d'un compte bancaire local~; une société étrangère ne pouvait ni ouvrir de compte ni retirer de won. \textbf{10.} Les bourses coréennes
n'autorisaient pas la vente à découvert, si bien que la prime ne pouvait être négociée qu'en faisant entrer des crypto-actifs. \textbf{11.} L'écart s'est élargi~: la demande locale a augmenté
plus vite que le capital d'arbitrage ne pouvait entrer par la jambe en monnaie fiduciaire. \textbf{12.} Les deux jambes étaient en crypto-actifs, qui circulaient
librement. \textbf{13.} Le fait que la contrainte soit restée active pendant des mois, à une échelle qu'aucun frais ne pouvait expliquer~: la limite était
une autorisation de déplacer de l'argent, et non une information. \textbf{14.} Une seule fois~: vendre de l'inventaire en Corée livre des won, qu'il faut encore
faire sortir pour racheter des crypto-actifs à l'étranger~; la contrainte active est la jambe en monnaie fiduciaire. \textbf{15.} Seulement si les \omterm{def:m3:blockchains-for-traders:stable}{stablecoins}
achetés en won pouvaient être vendus à l'étranger sans la même prime~; une prime à l'achat de \omterm{def:m3:blockchains-for-traders:stable}{stablecoins} en won est la même
contrainte sous une autre forme. \textbf{16.} C'est le prix d'une restriction légale, et non une incapacité à repérer un écart~; il
n'est une inefficience que par rapport à un monde sans restriction. \textbf{17.} Les résidents habilités selon les règles
et disposant de comptes à l'étranger, qui s'exposaient aux variations de prix pendant le transfert, aux limites légales et aux bourses détenant
leurs fonds. \textbf{18.} Que les prix se forment au sein de zones monétaires, et que les écarts entre elles mesurent le coût
de déplacement de l'argent, et non la qualité du prix. \textbf{19.} \emph{Résultat nommé~:} 4.03~\% avec un coût de rapatriement de 2~\%,
1.95~\% sans lui. \textbf{20.} Le temps et l'autorisation nécessaires pour transférer des actifs, et surtout de l'argent, entre
elles.
\end{solution}

\begin{solution}{iq:m3:spot-markets:1}
Un cycle passant par trois paires sur une même plateforme qui rend davantage de l'actif de départ qu'il n'en a dépensé. Il doit dépasser trois
frais preneurs et les fourchettes, et les teneurs de marché de la plateforme maintiennent les cours croisés dans cette bande, si bien que ce qui reste est réduit
et capté par les intervenants aux frais les plus faibles et les plus rapides.
\iqlookfor{la bande de frais et la concurrence.}
\end{solution}

\begin{solution}{iq:m3:spot-markets:2}
C'est un cours en USDT, dont la valeur en dollars est elle-même un taux de change négocié susceptible de bouger, peu en temps normal et beaucoup
en cas de \omterm{def:m3:blockchains-for-traders:stable}{décrochage} (depeg). Les soldes en \omterm{def:m3:blockchains-for-traders:stable}{stablecoins} sont évalués au prix de marché, et un arbitrage dollar--USDT est une position en USDT.
\iqlookfor{le cours croisé implicite et l'évaluation au prix de marché.}
\end{solution}

\begin{solution}{iq:m3:spot-markets:3}
Comparer ses données de transactions avec celles des plateformes réglementées~: premiers chiffres des tailles confrontés à la loi de Benford, concentrations aux
tailles rondes, la queue de la distribution des tailles~; comparer le volume à la \omterm{def:m3:spot-markets:depth}{profondeur affichée}, la co-variation des prix avec les principales plateformes,
et le trafic web ou les flux on-chain~; estimer la part de \omterm{def:m3:spot-markets:wash}{wash trading} à partir des arrondis par rapport à un indice de référence.
\iqlookfor{tests statistiques et recoupements avec des données indépendantes.}
\end{solution}

\begin{solution}{iq:m3:spot-markets:4}
Le coût de refermeture peut excéder l'écart~: frais, lenteur de crédit des dépôts ou des retraits, risque de prix lié au transfert,
limites aux mouvements de monnaie fiduciaire, accès aux comptes, ou doutes sur la solvabilité de la plateforme la plus chère (un écart peut être un spread de crédit).
\iqlookfor{latence, accès et risque de contrepartie.}
\end{solution}

\begin{solution}{iq:m3:spot-markets:5}
Risque de prix des crypto-actifs (couverture par contrats à terme), défaut de plateforme (limites par plateforme, balayage de l'excédent), \omterm{def:m3:blockchains-for-traders:stable}{décrochage} de \omterm{def:m3:blockchains-for-traders:stable}{stablecoin} (depeg)
(diversifier et évaluer au prix de marché), risque opérationnel lors des transferts (listes blanches d'adresses, rapprochements), et coût de rééquilibrage
lorsque les flux sont à sens unique (limites de déséquilibre, budget de rééquilibrage).
\iqlookfor{risque de marché, de crédit et opérationnel, chacun avec une limite.}
\end{solution}

\begin{solution}{iq:m3:spot-markets:6}
Estimer, par plateforme et par actif, le flux d'opportunités et leur dérive à sens unique d'après l'historique~; le coût et la durée de
déplacer l'inventaire entre chaque paire~; la limite de crédit de chaque plateforme. Résoudre périodiquement pour les inventaires cibles qui
maximisent l'edge capturé espéré diminué des coûts de rééquilibrage et de capital dans les limites (un problème de transport ou un petit
programme linéaire)~; rééquilibrer lorsque les écarts dépassent des seuils, en privilégiant les chaînes peu coûteuses et les heures creuses~; suivre les taux d'exécution et
recalibrer.
\iqlookfor{prévision des flux, des coûts, des limites, et un optimiseur à seuils.}
\end{solution}
